% Part: applied-modal-logics
% Chapter: temporal-logic
% Section: possible-histories

\documentclass[../../../include/open-logic-section]{subfiles}

\begin{document}

\olfileid{aml}{tl}{poss}

\olsection{التواريخ الممكنة}

النماذج العلائقية للمنطق الزمني التي استعملناها شديدة المرونة، لأننا لسنا
مضطرين إلى وضع أي قيد على علاقة النفاذ. وهذا يعني أن النماذج الزمنية يمكن
أن تتفرع في الماضي وفي المستقبل؛ لكن قد نرغب في تصور «جهوي» أكثر للتفرع،
ننظر فيه إلى متتاليات الأحداث بوصفها تواريخ ممكنة. ولا يستلزم ذلك بالضرورة
تغيير لغتنا، وإن كان بوسعنا أيضًا إضافة مؤثرينا الجهويين «العاديين»~$\Box$
و$\Diamond$، كما يمكننا التفكير في إضافة علاقات نفاذ معرفية لتمثيل تغيرات
معرفة الفاعلين عبر الزمن.

\begin{defn}
  \emph{نموذج التواريخ الممكنة} للغة الزمنية ثلاثي
  $\mModel{M}=\tuple{T,C,V}$، حيث:
  \begin{enumerate}
    \item $T$ مجموعة غير فارغة، تُفسر عناصرها بوصفها حالات في الزمن.
    \item $C$ مجموعة من المسارات الحسابية، أو \emph{التواريخ الممكنة}
      لنظام ما. وبعبارة أخرى،~$C$ مجموعة من المتتاليات~$\sigma$ للحالات
      $s_1$، $s_2$،~$s_3$، \dots، حيث كل~$s_i \in T$.
    \item $V$ دالة تسند إلى كل~!!{propositional
      variable}~$p$ مجموعة~$V(p)$ من النقاط الزمنية.
  \end{enumerate}
  ولتبسيط الأمور، سنفترض في العموم أيضًا أنه إذا كان تاريخ ما في~$C$، كانت
  جميع لواحقه فيها أيضًا. فمثلًا، إذا كانت~$s_1$، $s_2$،~$s_3$ متتالية
  في~$C$، كانت~$s_2$، $s_3$ وكذلك~$s_3$ فيها. وإذا ظهرت حالتان~$s_i$
  و$s_j$ في متتالية~$\sigma$، قلنا إن~$s_i \prec_\sigma s_j$ متى كان
  $i<j$. وحين يكون~$t \in V(p)$، نقول إن~$p$ \emph{صادق عند}~$t$.
\end{defn}

التغيير الوحيد ذو الصلة هو أننا حين نقوّم صدق !!a{formula} عند نقطة
زمنية~$t$ في نموذج~$\mModel M$، نفعل ذلك بالنسبة إلى تاريخ~$\sigma$
تظهر فيه~$t$ بوصفها حالة. ولا حاجة، مع ذلك، إلى تغيير دلالات أي من
ال!!{propositional variable}s أو الروابط الدالية الصدق؛ فهي جميعًا كما
كانت تمامًا في~\olref[sem]{defn:tmodels}، لأن أيًا منها لا يشير إلى
$\sigma$. لكننا نعيد الآن تعريف مؤثر المستقبل~$\Ftemp$، ونضيف مؤثرنا
$\Diamond$ بالنسبة إلى هذه التواريخ.

\begin{defn}\ollabel{defn:phmodels}
  \emph{صدق !!a{formula}~$!A$ عند~$t,\sigma$} في~$\mModel M$، ورمزه
  $\mSat{M}{!A}[t,\sigma]$:
  \begin{enumerate}
  \item\ollabel{defn:sub:phmodels-f} \indcase{!A}{\Ftemp
    !B}{$\mSat{M}{\indfrm}[t,\sigma]$ إذا وفقط إذا كان
    $\mSat{M}{!B}[t',\sigma]$ لبعض~$t' \in T$ بحيث
    $t \prec_\sigma t'$.}
  \item\ollabel{defn:sub:phmodels-diamond} \indcase{!A}{\Diamond
    !B}{$\mSat{M}{\indfrm}[t,\sigma]$ إذا وفقط إذا كان
    $\mSat{M}{!B}[t,\sigma']$ لبعض~$\sigma' \in C$ تظهر فيه~$t$.}
  \end{enumerate}
\end{defn}

ويمكن تعريف مؤثرات زمنية وجهوية أخرى بالطريقة نفسها. لكن بوسعنا الآن تمثيل
دعاوى تجمع بين الزمن والجهة. فمثلًا، يمكن أن نرمز إلى «لن يقع~$p$، لكنه كان
يمكن أن يقع» بالصيغة~$\lnot \Ftemp p \land \Diamond \Ftemp p$. وتصدق
هذه عند نقطة وتاريخ لا يصبح فيهما~$p$ صادقًا في حالة لاحقة، مع وجود تاريخ
بديل يصبح فيه~$p$ صادقًا في المستقبل.

\end{document}
